\chapter{Falsification of the Critiques}
\label{ch:critique-reassessment}

A critique is itself a claim. This chapter therefore subjects each material
objection from the first audit to the same rule applied to the source corpus:
name a possible falsifier, inspect the relevant repository surface, and retain
only the narrowest verdict supported by the result. A surviving critique is
not strengthened by rhetoric. A failed critique is narrowed or withdrawn.

The canonical record is
\path{data/registry/critique_evidence_ledger.json}. The table below is generated
from that ledger, and generation fails when an evidence path is missing. This
keeps the prose verdict, machine-readable status, and repository evidence on
one truth surface.

\small
\begin{longtable}{>{\raggedright\arraybackslash\bfseries}p{0.22\textwidth}
                  >{\raggedright\arraybackslash}p{0.43\textwidth}
                  >{\raggedright\arraybackslash}p{0.23\textwidth}}
  \caption{Critique-by-critique falsification record.}
  \label{tab:critique-reassessment}\\
  \toprule
  Critique & Falsification attempt and repository evidence & Verdict \\
  \midrule
  \endfirsthead
  \toprule
  Critique & Falsification attempt and repository evidence & Verdict \\
  \midrule
  \endhead
% Generated by scripts/generate_critique_evidence_table.py.
  The value c to the fourth divided by G proves force unification &
    Search the retained sources for field content, a common gauge structure, derived couplings, or a distinguishing prediction. The corpus supplies dimensional identities and interpretations but no such derivation. &
    Survives: The identity defines a force scale, not a unification mechanism. \\
  Cancellation of hbar is independent physical evidence &
    Expand Planck energy divided by Planck length. The cancellation is an algebraic consequence of the selected derived units. &
    Survives: No independent macroscopic quantum effect follows from the cancellation. \\
  E7 roots carry nontrivial P/Q parity &
    Classify all 126 generated roots. Every root belongs to the root lattice Q and therefore represents the zero coset. &
    Survives: The stated root-class map is trivial. \\
  No nontrivial Fourier-to-P/Q map can exist &
    Enumerate every homomorphism from Z\textasciicircum{}2 to the E7 quotient P/Q isomorphic to Z/2 and impose square-lattice D4 symmetry. The unique nontrivial invariant map is h(kx,ky)=(kx+ky) mod 2. &
    Falsified: A nontrivial symmetry-compatible map exists, but homomorphism additivity admits every exact triad and therefore supplies no nontrivial triad filter. \\
  A modular triad filter drives the simulation &
    Search the CPU, JAX, notebook, and high-resolution driver paths for quotient labels and a triad-admission predicate. Neither exists. &
    Survives: The proposed operator is absent from the implemented dynamics. \\
  The retained dynamics implement a beta plane &
    Trace every parameter into collision, forcing, and streaming. The notebook computes a beta value, but beta does not enter the evolution. &
    Survives: Arithmetic after the run is not beta-dependent dynamics. \\
  Root enumeration is nondeterministic &
    Repeat generation across hash seeds and inspect the set-to-array boundary. The tested float tuples were stable, but set iteration was an unstated dependency; the generator now sorts lexicographically. &
    Repaired: Enumeration is deterministic, while physical invariance under ordering remains unproved. \\
  Root order is physically irrelevant &
    Compare the canonical density scaffold with controlled permutations. Reversal is exactly invariant because opposite roots share absolute coordinates, while a one-place cyclic shift changes the perturbation by 0.131 in relative L2. &
    Falsified: One symmetry-related permutation is harmless, but index weighting remains load-bearing. \\
  The high-resolution run has catastrophic mass growth &
    Read all retained Parquet states. Relative drift reaches only 5.60e-6 at step 500. &
    Narrowed: Catastrophic growth belonged to the separate historical CPU demonstration. \\
  The historical CPU demonstration conserves mass &
    Sum the historical D2Q9 equilibrium at nonzero velocity. Its coefficients applied the sound-speed factors twice, producing density times one plus nine times speed squared. &
    Repaired: The historical critique survives, and the current equation and validator now pass conservation tests. \\
  The retained run supports four zonal jets &
    Measure the final state with a fixed diagnostic. Mode 17 dominates the small zonal-mean profile, which is only 1.22e-3 of total velocity RMS; the vorticity moment ratio is 0.992. &
    Survives: A peak index alone does not establish a strong zonal regime. \\
  The reported Rhines comparison is independent &
    Trace beta\_eff in the notebook. It is reconstructed from the observed velocity and peak wavenumber used in the comparison. &
    Survives: The reported agreement is circular. \\
  Affine symmetry transfers to the fluid through the E8 label &
    Search for currents, operator products, a stress tensor, and Ward tests. The repository supplies none. &
    Survives: The Sugawara result does not transfer to the fluid model by naming. \\
  Imaginary algebraic roots are evanescent spatial modes &
    Search for a dimensional map between the Kac-Moody bilinear form and a spatial dispersion relation. None is defined. &
    Survives: The identification conflates different quadratic forms. \\
  The corpus contains no empirical materials evidence &
    Inspect the tourmaline and zero-point-vibration sources. They include ordinary electrical and thermal measurements and standard lattice-vibration calculations. &
    Falsified: Relevant conventional evidence is present. \\
  The corpus supports novel tourmaline or excess-energy effects &
    Search for calibrated raw data, complete energy balances, blinded controls, and independent replication of the proposed enhancement. &
    Survives: Conventional evidence does not test the novel coupling. \\
  Two indexed JSON artifacts are unusable &
    Parse the historical files, identify truncation, and rerun deterministic generators with schema-valid output and invariant tests. &
    Repaired: The historical critique survives; the current artifacts are parseable and regenerated. \\
  Registry inclusion establishes scientific validity &
    Compare identity checks with parsing, schema, conservation, and model gates. Registry identity is necessary but not sufficient. &
    Survives: A digest establishes byte identity, not scientific validity. \\
  The retained fractal dimensions are converged &
    Compare retained estimates with exact controls and inspect scale-window stability. Multiple estimates disagree and no refinement study closes the gap. &
    Survives: The values are finite-resolution diagnostics. \\
  Layered harmonics measure an E7 physical effect &
    Trace the retained schedule to a physical observable or discriminating control. It is a parameter construction without such a link. &
    Survives: The artifact records a chosen construction rather than a measured E7 effect. \\
  A matched dynamical ablation is the uniquely smallest next test &
    Separate zero-cost structural checks from dynamical experiments. Exhaustive triad inventory can reject the proposed filter before any simulation. &
    Falsified: The exhaustive audit closes the homomorphic filter before dynamics; a beta-plane run now serves as a negative matched control rather than a rescue test. \\
  \bottomrule

\end{longtable}
\normalsize

The reassessment changes the paper in three ways. It withdraws absolute claims
about the absence of conventional materials evidence, confines catastrophic
mass growth to the defective historical CPU run, and records repaired artifacts
as current regressions rather than continuing to call them malformed. The main
physical objections survive because the missing maps, operators, and controls
remain missing after those corrections.
